% 整数（综述）
% license CCBYSA3
% type Wiki

本文根据 CC-BY-SA 协议转载翻译自维基百科\href{https://en.wikipedia.org/wiki/Integer}{相关文章}

整数（integer）是指 **零（0）**、**正自然数（1, 2, 3, …）**，或 **正自然数的相反数（−1, −2, −3, …）**。[1]正自然数的相反数（或加法逆元）称为 **负整数（negative integers）**。[2]所有整数构成的集合通常记作 **粗体 Z** 或 **黑板粗体**：$\mathbb{Z}$.[3][4]  

自然数集合  $\mathbb{N}$是整数集合$\mathbb{Z}$的子集，而 $\mathbb{Z}$ 又是有理数集合$\mathbb{Q}$的子集，而 $\mathbb{Q}$ 本身是实数集合$\mathbb{R}$的子集。[a]与自然数集合类似，整数集合$\mathbb{Z}$是可数无限集。
整数可以被视为一种实数，其表示中不含分数部分。
例如：$21, 4, 0, -2048$ 都是整数；而 $9.75, 5+\tfrac{1}{2}, \tfrac{5}{4}, \sqrt{2}$ 都不是整数。\[7]

整数构成了包含自然数的最小的 **群（group）** 与最小的 **环（ring）**。
在代数数论中，整数有时被称为 **有理整数（rational integers）**，以区别于更一般的 **代数整数（algebraic integers）**。
实际上，（有理）整数就是那些既是代数整数、又是有理数的数。

---

要不要我帮你把 **整数在集合包含关系中的层级结构** 画成一张 **集合嵌套图**（$\mathbb{N} \subset \mathbb{Z} \subset \mathbb{Q} \subset \mathbb{R}$）？
